% Figure: explicit vs implicit accuracy/cost trade-off for the heat
% equation. Top panel: midpoint decay of sin(pi x) under FTCS,
% BTCS, Crank-Nicolson and exact, on a coarse time step.
% Bottom panel: stable time-step ceiling vs grid resolution for
% explicit (quadratic) and implicit (accuracy-limited) schemes.
\begin{figure}[htbp]
\centering

% --- Top panel: midpoint trajectory u(1/2, t) ---
\begin{tikzpicture}
\begin{axis}[
  width=12cm, height=5cm,
  title={\small (a) midpoint decay $u(1/2,t)$, coarse step $r = 5$},
  xlabel={time $t$},
  ylabel={$u(1/2,t)$},
  xmin=0, xmax=0.25,
  ymin=0, ymax=1.05,
  grid=both,
  major grid style={engblack!20},
  legend pos=north east,
  legend cell align=left,
  font=\scriptsize,
  samples=120, domain=0:0.25,
]
% Exact decay e^{-pi^2 t}.
\addplot[engblack, very thick] {exp(-pi^2 * x)};
\addlegendentry{exact $e^{-\pi^2 t}$}
% BTCS one-step factor per step of size dt; with r=5 dx=1/20 dt=5/400=0.0125.
% Approximate implicit-Euler decay (1/(1+pi^2 dt))^{t/dt}: over-damps slightly.
\addplot[engred, very thick, mark=square*, mark size=1.2pt, samples=21, domain=0:0.25] {
  (1/(1+pi^2*0.0125))^(x/0.0125)
};
\addlegendentry{BTCS (first order, over-damps)}
% Crank-Nicolson factor ((1-pi^2 dt/2)/(1+pi^2 dt/2))^{t/dt}: closer.
\addplot[engblue, very thick, mark=*, mark size=1.2pt, samples=21, domain=0:0.25] {
  ((1-pi^2*0.0125/2)/(1+pi^2*0.0125/2))^(x/0.0125)
};
\addlegendentry{Crank-Nicolson (second order)}
\end{axis}
\end{tikzpicture}

\vspace{1.5ex}

% --- Bottom panel: stable dt ceiling vs N ---
\begin{tikzpicture}
\begin{loglogaxis}[
  width=12cm, height=5cm,
  title={\small (b) usable time step vs grid points $N$},
  xlabel={grid points $N$ (with $\Delta x = 1/N$)},
  ylabel={$\Delta t$},
  xmin=10, xmax=1000,
  grid=both,
  major grid style={engblack!20},
  legend pos=south west,
  legend cell align=left,
  font=\scriptsize,
]
% Explicit ceiling: dt = dx^2/(2 alpha), alpha=1 -> dt = 1/(2 N^2).
\addplot[engred, very thick, domain=10:1000, samples=40] {1/(2*x^2)};
\addlegendentry{explicit CFL ceiling $\Delta t = \Delta x^{2}/2$}
% Implicit accuracy-limited: dt ~ const (say accuracy target dt=0.01).
\addplot[engblue, very thick, domain=10:1000, samples=2] {0.01};
\addlegendentry{implicit accuracy budget $\Delta t \approx$ const}
\end{loglogaxis}
\end{tikzpicture}

\caption[Explicit versus implicit trade-off]{The working trade-off
between explicit and implicit time-stepping for the heat equation
$u_t = u_{xx}$. Panel (a): the midpoint decay of the fundamental
mode $\sin(\pi x)$ under a deliberately coarse step ($r = 5$, far
above the explicit CFL bound, so FTCS would diverge and is not
shown). BTCS is stable but over-damps because its one-step factor
$1/(1 + \pi^2 \Delta t)$ underestimates $e^{-\pi^2 \Delta t}$ at
first order. Crank-Nicolson tracks the exact decay closely because
its factor $(1 - \pi^2 \Delta t/2)/(1 + \pi^2 \Delta t/2)$ is the
second-order Pad\'e approximant. Panel (b): the usable time step
against grid resolution. The explicit CFL ceiling falls as $\Delta
x^{2}$ (slope $-2$ on log-log axes), so doubling resolution quarters
the step. The implicit step is limited only by accuracy and stays
roughly constant, so on fine grids the implicit scheme takes orders
of magnitude fewer steps. The crossover is where per-step cost of
the tridiagonal solve no longer pays for the larger step it
permits.}
\label{fig:vol02:ch12:explicit-implicit}
\end{figure}
